% SPDX-License-Identifier: Apache-2.0
\section{The Other Side of a Bundle, Derived}
\label{sec:x-hostgen}

\code{link::<B>()} makes both sides of a bundle in one line of a unit
of units, as Section~\ref{sec:x-link} does, for a \code{B} that
implements \code{Link}. \code{Link} was written by hand, and so was the
struct the other side holds. \code{\#[derive(Link)]} writes both
(issue 1695). Beside \code{Lane} it writes \code{LaneHost}, each field
the end opposite \code{Lane}'s, an \code{Rx} a \code{Tx} and an
\code{Out} an \code{In}, with \code{Ports} on it and the \code{Link}
between the two. An array of lanes and a nested bundle turn round
too, so \code{Lanes} holding \code{[Lane; 2]} gets \code{LanesHost}
holding \code{[LaneHost; 2]}. A struct generic over its side needs no
derive of its own: Section~\ref{sec:x-flipped}'s \code{\#[derive(Ports)]}
writes its \code{Link}.

Here \code{Feed} takes \code{LanesHost} and sends a count on both
lanes, the second a hundred above the first, and counts the first
lane's acknowledgements. \code{Take} takes \code{Lanes}, keeps each
lane's last beat and acknowledges it. \code{Pair} joins them with
\code{link::<Lanes>()}. Both ways are channels, so the order the two
units run in does not matter. The run checks every cycle from the
third, and the netlist is checked against the run under both
simulators.

\begin{figure}[htbp]
\centering
\input{hostgen_timing.tex}
\caption{The count sent, the last beat of each lane, and the first
lane's acknowledgements.}
\label{fig:x-hostgen}
\end{figure}

\lstinputlisting[caption={\code{ex\_hostgen.rs}. Run with \code{bazel run //lib/examples:ex\_hostgen}.},label={lst:x-hostgen}]{lib/examples/ex_hostgen.rs}

\lstinputlisting[style=ascii,caption={What \code{ex\_hostgen} printed when the build ran it: each cycle's acknowledgements and lanes, then the lowering.},label={lst:x-hostgen-out}]{out_hostgen.txt}
